\section{Discussion}
\label{sec:discussion}
\subsection{Interpretation}
The route study supplies the current-version evidence. Registered tasks retain
task, approval, receipt, and sample associations; direct REST retains
operation, readback, and sample records without a task edge. Each route used
the scene's bound DAQ and DCW nodes, so binding is observed as the precondition
under which the recorded reads and writes succeeded. The two simulator scenes
use different signal and protocol mappings, and the fixed-call and CLI routes
separate task-path behavior from an external model decision. Similar
post-trial samples are consistent with a shared handler; the annealing
difference is consistent with the 735 versus 745\,$^{\circ}$C setpoints. But
separate environments and one sample per task or call do not establish route
equivalence.

The approval-context test isolates the tuple comparison under mock state:
removing that comparison admitted product/run changes that the retained checks
alone did not reject, while line and recipe changes remained blocked. The
blueprint replay shows that one documented coating line can be instantiated
from configuration objects with four DCW setpoints, seven DAQ nodes, and four
MES forms, and that the verifier can check sample availability and write
behavior without altering the protected demonstration line.

The four archived missions apply one operation--response--judgment workflow to
distinct objectives: paired holding adjustments for molding, three-control
thickness correction for BOPET, throughput under a hardness constraint for
annealing, and staged aeration and dosing for wastewater. Most requested
changes are followed by a recorded process observation; reverted candidates and
the 101.3-HV event-text value are marked in the figure. The missions show how a
common operation path accommodates different actuator structures and target
types. They are historical scripted direct-node records, not current-route
trials, and each is a single trajectory.

The process results bound the interpretation. Cast-film samples were
1.31--2.31\,$\mu$m below the 48\,$\mu$m lower bound, and annealing samples were
20.40--25.60\,$^{\circ}$C below 700\,$^{\circ}$C. Because five of six tasks
started below their band, the post-trial shortfall is not evidence that the
selected setpoints caused it; the one above-band start cannot be resolved with
a single 20-s probe. The archived missions end inside their bands, but they
measure different process variables and are not a controlled comparison with
the route study.

Taken together, the evidence supports a software separation between task
reasoning, operation resolution, scene mapping, and process observation; the
archived missions show that this sequence survives changes in actuator
structure and target type.

\subsection{Scope and Limitations}
Only the reported adapters and scenes were exercised. The route runs and the
C01 replay used working-tree builds rather than pinned releases; the route
study used the SQLite fallback because the production MQTT/Timescale backends
were not used; and the approval evidence is mock state with synthetic
approvals. The four archived missions are one seed-42 trajectory per
scenario, generated through an earlier scripted direct-node interface, and are
not current-route, repeatability, or physical-plant evidence. The study also
does not establish adapter interchangeability, configuration reuse, normal
online DAQ behavior, or a platform baseline.
The feature comparison positions the implementation question; it does not show
that AgentWorkShop is more efficient or more capable than another platform. A
comparative study would need a shared task set, explicit integration-effort or
record-completeness measures, and matched deployments. Operation restrictions
do not establish security
\cite{greshake2023injection,liu2024pinjection,iec62443-3-3} or replace
independent PLC/safety-instrumented system (SIS) safeguards \cite{iec61511}.
Synthetic approval also does not measure human-review performance
\cite{parasuramanmanzey2010complacency,bainbridge1983ironies}.
Scope and channel-policy refusals are not uniformly persisted; auditability is
limited to the records and denial classes described in Section~IV.

\section{Conclusion}
\label{sec:conclusion}
AgentWorkShop connects industrial task workflows to scene-specific operations
through shared dispatch, configured drivers, and service-owned process records.
A route study across four entry paths preserved entry-specific records and used
scene-bound DAQ and DCW nodes; a mock-state approval-context test rejected
stale approval tuples before dispatch; a blueprint replay instantiated one
documented coating-line configuration in simulation; and four archived
simulator missions illustrated the operation--response--judgment sequence
across distinct process objectives. The study establishes record linkage
across the tested entry paths, stale-context rejection in mock state, and one
simulated onboarding instance; it does not establish process improvement,
configuration reuse, general portability, or comparative platform benefit.
